%\documentclass{beamer}
\documentclass[aspectratio=169]{beamer}
\usetheme{Madrid}
\usecolortheme{default}

% Packages
\usepackage[utf8]{inputenc}
\usepackage{graphicx}
\usepackage{booktabs}
\usepackage{algorithm}
\usepackage{algorithmic}
\usepackage{multimedia} 

% Title Information
\title[Lab 06: Cloud ML]{Lab Exercise 06: Machine Learning in the Cloud with Microsoft Azure}
\subtitle{ Human-like Shogi AI}
\author{Takanori Nagashima}
\institute{Introduction to Cloud Computing}
\date{\today}
\AtBeginSection[]{
  \begin{frame}{Outline}
    \tableofcontents[currentsection]
  \end{frame}
}

\begin{document}

% --- Slide 1: Title ---
\frame{\titlepage}

\begin{frame}{Outline}             % \begin{frame}{Title} sets the slide heading
    \tableofcontents               % Generates TOC (auto-linked)
\end{frame}

\section{Introduction}
% --- Slide 2: ---
\begin{frame}{Project Overview: Human-like Shogi AI}
    \begin{block}{Concept}
        Development of a Shogi AI that emulates human-like decision-making processes rather than purely focusing on win-rate maximization.
    \end{block}
    
    \begin{itemize}
        \item \textbf{Data Source:} Supervised learning based on professional game records (\textit{.kif}).
        \item \textbf{Purpose:} Create a human-like tutor to counter the issue of current AIs being too strong.
        \item \textbf{Platform:} Utilizing \textbf{Microsoft Azure} for large-scale model training and deployment.
    \end{itemize}
\end{frame}


% --- Slide 2: What is Shogi? ---
\begin{frame}{What is Shogi?}
    \begin{center}
        \includegraphics[width=0.3\textwidth]{assets/board_image.jpeg}
    \end{center}
    \begin{itemize}
        \item A traditional Japanese strategy board game.
        \item Similar to Chess, but featuring a unique "Drop Rule" where captured pieces can be returned to the board as your own.
        \item Known for its extreme complexity ($>$ Chess) and deep strategic possibilities.
    \end{itemize}
\end{frame}


% --- Slide 3: Azure Setup ---
\begin{frame}{Azure Setup}
    \begin{itemize}
        \item \textbf{Resource Provisioning:} Used \textbf{Azure CLI} for environment setup.
        \item \textbf{Deployment Script:} \texttt{01-azure.sh}
            \begin{itemize}
                \item This script automates the creation and configuration of the necessary Azure resources.
            \end{itemize}
        \item \textbf{Region Selection: \texttt{japaneast}}
            \begin{itemize}
                \item Chosen for its geographical proximity to the target Japanese user base.
                \item Intended to optimize network performance and minimize latency.
            \end{itemize}
    \end{itemize}
\end{frame}

\section{Data Collection}
% --- Slide 3: Data Collection - 1 ---
\begin{frame}{Data Collection - 1}
    \textbf{Target Platform:} Shogi Wars (The most popular Shogi application)
    \begin{itemize}
        \item URL: \url{https://shogiwars.heroz.jp/}
    \end{itemize}

    \begin{center}
        \includegraphics[height=0.3\textheight]{assets/shogiwars_image.jpeg}
    \end{center}

    \begin{block}{The Problem}
        The Shogi Wars provider officially offers neither game record (kif) data for public download nor a public API interface.
    \end{block}
\end{frame}

% --- Slide 4: Data Collection - 2 ---
\begin{frame}{Data Collection - 2: Technical Approach}
    \begin{itemize}
        \item \textbf{Exploration of Methods:} Two potential approaches:
            \begin{enumerate}
                \item Investigating the request structure to identify data endpoints.
                \item \textbf{Extracting KIF strings directly from HTML source.}
            \end{enumerate}
        \item \textbf{Investigation via \texttt{curl}:} 
            \begin{itemize}
                \item Manually tested multiple pages using \texttt{curl} scripts.
                \item Confirmed that no specific data endpoints exist; the kif data was embedded within the HTML.
            \end{itemize}
        \item \textbf{Final Decision:} Selected the method of scraping KIF strings from HTML using Python.
    \end{itemize}
\end{frame}

% --- Slide 5: Data Collection Methodology ---
\begin{frame}{Data Collection - 3: Basic Idea}
    \begin{algorithm}[H]
        \caption{Data Extraction and Conversion Pipeline}
        \begin{algorithmic}[1]
            \STATE \textbf{Identify Users:} Access the "Yearly Championship" page to extract active \texttt{user\_id}s.
            \STATE \textbf{Filter Games:} For each \texttt{user\_id}, extract \texttt{game\_id}s filtered by game type: \textbf{sb (3-minute games)}.
            \STATE \textbf{Storage:} Store the \texttt{user\_id} and \texttt{game\_id} mappings in a JSON file.
            \STATE \textbf{Scrape KIF:} Download and extract raw \texttt{.kif} strings from the associated HTML pages.
            \STATE \textbf{Format Conversion:} 
            \begin{itemize}
                \item Convert \texttt{.kif} $\rightarrow$ \texttt{.csa} format.
                \item Tool: \textbf{ShogiHome} (\url{https://sunfish-shogi.github.io/shogihome/})
            \end{itemize}
        \end{algorithmic}
    \end{algorithm}
\end{frame}

\begin{frame}{Data Collection - 4: Extracting Session Data}
    \begin{itemize}
        \item \textbf{Challenge:} Accessing protected game records requires valid authentication tokens.
        \item \textbf{Investigation:} Identifying \texttt{Websession} values and \texttt{Cookie} strings through browser developer tools and manual \texttt{curl} tests.
    \end{itemize}

    \begin{center}
        \includegraphics[width=0.4\textwidth]{assets/websession.png}
    \end{center}
\end{frame}

\begin{frame}{Data Collection - 5: Main Code}

    \begin{block}{Solution: Data Scraping}
        To collect the necessary dataset, the following script was utilized to clone game records:
        \begin{itemize}
            \item \url{https://github.com/LoveKapibarasan/kifs}
        \end{itemize}
    \end{block}

    \begin{alertblock}{Important Remark: Ethical Scraping}
        \begin{itemize}
            \item \textbf{Do not overload the server:} Implemented appropriate intervals (sleep timers) between requests.
            \item Respecting the target platform's stability is crucial when collecting large-scale datasets.
        \end{itemize}
    \end{alertblock}
    
\end{frame}

\begin{frame}{Data Collection - 6: Azure Infrastructure}
    \begin{itemize}
        \item \textbf{Computational Resource: Azure Virtual Machines}
        \begin{itemize}
            \item Used for running the Python scraping scripts and data conversion.
            \item \textbf{Spot Instance:} Utilized Azure Spot Instances to significantly reduce operational costs (up to 90\% discount) for non-urgent background tasks.
        \end{itemize}
        
        \item \textbf{Data Persistence: Azure Blob Storage}
        \begin{itemize}
            \item Collected \texttt{.kif} and \texttt{.csa} files are directly uploaded to Blob Storage.
            \item Ensures data durability and accessibility for the subsequent Machine Learning phase.
        \end{itemize}
    \end{itemize}

    \begin{block}{Why Spot Instances?}
        Since data collection is a separable and resumable task, Spot Instances provided the ideal balance between high performance and low cost, even with the risk of eviction.
    \end{block}
\end{frame} 


\begin{frame}{Azure Setup: (02--04)}
    \begin{itemize}
        \item \textbf{02-azure.sh}
            \begin{itemize}
                \item Initiates the deployment environment and run the scripts.
            \end{itemize}
        \item \textbf{03-azure.sh}
            \begin{itemize}
                \item Provisions the Virtual Machine with Ubuntu Server 24.
            \end{itemize}
        \item \textbf{04-azure.sh}
            \begin{itemize}
                \item Automates the creation of the Azure Blob Storage container.
                \item Uploads the collected and converted Shogi game records (\texttt{.kif} / \texttt{.csa}) to the cloud for persistence.
            \end{itemize}
    \end{itemize}
\end{frame}

\begin{frame}{First Iteration Results}
    \begin{columns}
        \begin{column}{0.5\textwidth}
            \begin{itemize}
                \item \textbf{Data:} $\approx$ 6,000 kifs
                \item \textbf{Training Time:} 48 min
                \item \textbf{Status:} Pipeline validated
            \end{itemize}
        \end{column}
        \begin{column}{0.5\textwidth}
            \includegraphics[width=\textwidth]{assets/kif_to_csa.png}
        \end{column}
    \end{columns}
\end{frame}

\section{Training}

\begin{frame}{Development Environment \& Training}
    \begin{itemize}
        \item \textbf{Infrastructure:} 
            \begin{itemize}
                \item \textbf{Azure Machine Learning:} Managed environment for final integration and scalable deployment experiments.
                \item \textbf{Google Colab:} Primary training environment (Selected for cost-efficiency and resource optimization\footnote{Major training phases were migrated to Colab to manage budget constraints due to high Azure compute costs.}).
            \end{itemize}
        \item \textbf{Model Architecture:} Deep Reinforcement Learning with Shogi-specific residual blocks.
        \item \textbf{Style:} Human-like \textbf{Furibisha} (Ranging Rook).
        \item \textbf{Reference:} \texttt{05-azure.sh} (Provisioning script)
    \end{itemize}
\end{frame}

\section{Evaluation}

\begin{frame}{Result: Practical Performance}
    \begin{columns}[c]
        \begin{column}{0.45\textwidth}
            \begin{block}{AI Evaluation}
                \begin{itemize}
                    \item \textbf{Strength:} $\approx$ 1-kyu
                    \item \textbf{Style:} \textbf{Furibisha} (Ranging Rook)
                    \item \textbf{Strategy:} Professional human-like move selection.
                \end{itemize}
            \end{block}
            \small
            \textbf{Human-likeness:} Unlike typical engines, this model prefers aesthetic and tactical flexibility.
        \end{column}

        \begin{column}{0.5\textwidth}
            \begin{center}
                \begin{figure}
                    \href{https://vimeo.com/1186695731}{
                        \includegraphics[width=\textwidth]{assets/video_thumbnail.png}
                    }
                    \caption{\footnotesize \textcolor{blue}{\underline{Click to Watch Match Video (Vimeo)}}}
                \end{figure}
                \scriptsize ※ External browser will open for playback.
            \end{center}
        \end{column}
    \end{columns}
\end{frame}

\section{Deployment \& Architecture}


\begin{frame}{Communication \& Protocol}
    \begin{itemize}
        \item \textbf{Server Implementation:} 
            \begin{itemize}
                \item Utilized \texttt{shogi-server} for CSA protocol matching.
                \item Configured to listen on \textbf{Port 4081} (CSA Standard Port).
            \end{itemize}
        
        \item \textbf{CSA Protocol (ver. 1.2.1):}
            \begin{itemize}
                \item \textbf{Model:} TCP/IP server-client model.
                \item \textbf{Standard Port:} 4081.
                \item \textbf{Format:} Line-based 7-bit ASCII strings (LF terminated).
                \item \textbf{Reference:} \url{http://www2.computer-shogi.org/protocol/}
            \end{itemize}
    \end{itemize}

    \begin{block}{Core Mechanism}
        Commands like \texttt{LOGIN}, \texttt{AGREE}, and moves (e.g., \texttt{+7776FU}) allow standardized interaction between various engines and the shogi-server.
    \end{block}
\end{frame}


\begin{frame}{Cloud-Native Implementation with Azure \& Docker}
    \begin{itemize}
        \item \textbf{Core Engine:} \texttt{shogi-server}\footnote{https://github.com/shogi-server/shogi-server}
            \begin{itemize}
                \item Implemented in \textbf{Ruby 3.2+}.
                \item Compliant with \textbf{CSA Server Protocol Ver 1.2.1}.
            \end{itemize}
        
        \item \textbf{Infrastructure (Azure-Native):}
            \begin{itemize}
                \item \textbf{Azure Container Registry (ACR):} Managed private Docker images and facilitated automated deployment.
                \item \textbf{Azure Container Instances (ACI):} Leveraged for rapid, serverless container execution, exposing standard \textbf{Port 4081} to the public internet.
                \item \textbf{Reference:} \texttt{06-azure.sh} (Provisioning script)
            \end{itemize}
    \end{itemize}
\end{frame}

\section{How to Play}

\section{Web UI Deployment}

\begin{frame}{Web UI Deployment: Shogihome}
    \begin{columns}[c]
        \begin{column}{0.45\textwidth}
            \begin{itemize}
                \item \textbf{Automation Script:} \texttt{07-azure.sh}
                \item \textbf{Host Platform:} Azure App Service (Node.js 22-lts)
                \item \textbf{Outcome:} Frontend interface successfully bundled and deployed to the cloud infrastructure.
            \end{itemize}
        \end{column}

        \begin{column}{0.55\textwidth}
            \begin{center}
                \includegraphics[width=0.95\textwidth]{assets/terminal.png}
            \end{center}
        \end{column}
    \end{columns}
\end{frame}

\begin{frame}{How to Play - 1: Connect to the Server}
    \begin{columns}[c]
        \begin{column}{0.5\textwidth}
            \begin{block}{Scan to Access}
                Scan the QR code to access the Shogi client platform or view deployment details.
            \end{block}
            \begin{itemize}
                \item \textbf{Server Address:} \texttt{lovekapibarasan.org}                \item \textbf{Port:} \texttt{4081} (CSA Standard)
                \item \textbf{Engine Name:} \texttt{Server Engine}
            \end{itemize}
            \vspace{0.3cm}
        \end{column}

        \begin{column}{0.5\textwidth}
            \begin{center}
                \includegraphics[width=0.6\textwidth]{assets/qrcode.png}
                \vspace{0.2cm}
                \footnotesize \url{https://shogihome.lovekapibarasan.org}
            \end{center}
        \end{column}
    \end{columns}
\end{frame}

\begin{frame}{How to Play - 2: Language Settings}
    \begin{columns}[c]
        \begin{column}{0.45\textwidth}
            \begin{itemize}
                \item \textbf{Language Selection:} Navigate to the settings menu to toggle between Japanese and English.
            \end{itemize}
        \end{column}

        \begin{column}{0.55\textwidth}
            \begin{center}
                \includegraphics[width=0.9\textwidth]{assets/lang.png}
            \end{center}
        \end{column}
    \end{columns}
\end{frame}

\begin{frame}{How to Play - 3: Select the Bot Engine}
    \begin{columns}[c]
        \begin{column}{0.5\textwidth}
            \begin{block}{Steps to Start a Game}
                Follow these simple steps to challenge the human-like AI:
            \end{block}
            \begin{enumerate}
                \item Open the Web UI panel in your browser and select "Local Match".
            \end{enumerate}
        \end{column}

        \begin{column}{0.5\textwidth}
            \begin{center}
                \includegraphics[width=0.9\textwidth]{assets/how-to-play-1.png}
            \end{center}
        \end{column}
    \end{columns}
\end{frame}

\begin{frame}{How to Play - 4: Matchmaking with the Bot}
    \begin{columns}[c]
        \begin{column}{0.5\textwidth}
            \begin{block}{Steps to Start a Game}
                Follow these simple steps to challenge the AI:
            \end{block}
            \begin{enumerate}
                \setcounter{enumi}{1}
                \item Select \textbf{\texttt{Server Engine}} from the engine list.
                \item Send a challenge request to begin the match directly on the web interface.
            \end{enumerate}
        \end{column}

        \begin{column}{0.5\textwidth}
            \begin{center}
                \includegraphics[width=0.9\textwidth]{assets/how-to-play-2.png}
            \end{center}
        \end{column}
    \end{columns}
\end{frame}



\begin{frame}{Networking \& Infrastructure Trial}
    \begin{itemize}
        \item \textbf{Short-term Trials:} Evaluated various cloud AI services and managed networking components over a short period.
        \item \textbf{Architecture Pivot:} Attempted an \textbf{Azure Bastion} setup to bridge local GPU resources to the cloud, but abandoned it due to excessive pricing.
        \item \textbf{Final Solution:} Cost-effective hybrid setup using \textbf{WireGuard} and \textbf{Cloudflare}.
    \end{itemize}
\end{frame}

\section{Summary}


\begin{frame}{Acknowledgments}
    \begin{block}{Acknowledgments}
        Special thanks to the open-source foundations and projects that made this entire environment possible:
        \vspace{0.5em}
        \begin{itemize}
            \item \textbf{Shogihome:} For providing the comprehensive and intuitive frontend web interface.
            \item \textbf{shogi-server:} For managing the core CSA engine matchmaking system and protocol handling.
        \end{itemize}
    \end{block}
\end{frame}

\end{document}
